% Template; to be used with:
%          spconf.sty  - ICASSP/ICIP LaTeX style file, and
%          IEEEbib.bst - IEEE bibliography style file.
% --------------------------------------------------------------------------
\documentclass{article}
\usepackage{spconf,amsmath,graphicx}

% Example definitions.
% --------------------
\def\x{{\mathbf x}}
\def\L{{\cal L}}

% Title.
% ------
\title{Your project title}
%
% Single address.
% ---------------
\name{First1 Last1,  First2 Last2, etc}
\address{NUS-ISS, National University of Singapore, Singapore 119615}

\begin{document}
%\ninept
%
\maketitle
%

\begin{abstract}

The abstract should consist of one paragraph describing the motivation for your project and a high-level explanation of the methodology you used and the results obtained. Note: this \textit{two-column} project report template is modified from that used in Stanford University CS230, https://cs230.stanford.edu/. 

\end{abstract}
%

\textbf{Keywords}: Please provide (up to four) keywords

\section{Introduction}
\label{sec:intro}

Explain the problem statement and why it is important. Add one paragraph to clearly state the highlights (i.e., the selling point) of your project.

\section{Literature review}
Present relevant references (e.g., papers, industrial products), and discuss their strengths and weaknesses.

\section{Proposed approach}
Describe your proposed system. Use a system architecture or flowchart to illustrate your proposed system. For each module, give a detailed description of how it works.

\section{Experimental results}

\subsection{Dataset}
Describe your dataset. Include examples of your data in the report (e.g., include an image, show a waveform, etc.).

\subsection{Implementation details}

Data augmentation? Model training configurations, e.g., optimizer, epoch, batch size, learning rate, and GPU specification.

\subsection{Performance metrics}
What metrics are used: accuracy, precision, etc? Provide equations for the metrics.

\subsection{Experimental results}
You need to have both \textit{quantitative} and \textit{qualitative} results.

\subsection{Ablation study}
Conduct an ablation study to evaluate how various components contribute to its final performance (Table \ref{table1}).

\begin{table}[h]
\caption{The ablation study.}\label{table1} \centerline{
    \begin{tabular}{cc|cc}
    \hline\hline
    Module A & Module B & Metric X & Metric Y\\\hline
    - & $\surd$ & XXX & XXX \\
    $\surd$ & - & XXX & XXX \\
    $\surd$ & $\surd$ & XXX & XXX \\\hline
    \end{tabular}
    }
\end{table}

\subsection{Discussions and limitations}
Include some examples of where your approach failed and a discussion of why certain approaches failed or succeeded. 


\section{Conclusions and future work}
Summarize your report. For future work, if you had more time, more team members, or more computational resources, what would you explore? It would be good to have around 10-15 references \cite{He16CVPR} for the whole report.

\section{Author contributions}
Describe the contributions of each team member.

\section{AI Tool Declaration}

[Sample declaration copied from Canvas]

The authors didn't use any AI tool. The authors used [AI tool name, e.g., GPT-5.1] to [describe specific uses: e.g., generate ideas, format paragraphs, improve expression, analyse effectiveness, create images and illustrations, produce drafts, refine, and/or finalise my assignment]. The authors are responsible for the content and quality of the submitted work.

\bibliographystyle{IEEEbib}
\bibliography{references}

\end{document}
